\begin{frame}{One set of settings}
\vspace{.8em}
\centering
{\Large Integral threshold: \textbf{\textcolor{pennred}{0.3}}}\par
\vspace{1.8em}
{\Large Vertex finding: \textbf{\textcolor{pennred}{$|\Delta z|\leq0.1$ mm}}}\par
\vspace{1.8em}
{\Large Vertex categories: \textbf{\textcolor{pennred}{ATLAS-inspired diagnostic}}}\par
\vspace{.45em}
{\small Clean / Matched, Merged, Split, Fake}\par
\vfill
\sourcecaption{0.5 is a cross-check only. The all-vertex position metric is our diagnostic; the ITk paper uses 0.1 mm for HS finding. Current categories use unit track weights and an unlabelled-track proxy for Fake.}
\end{frame}

\begin{frame}{HL-LHC: vertex separation at integral cuts 0.3 and 0.5}
\centering
\includegraphics[width=.99\linewidth,height=.78\textheight,keepaspectratio]{figures/resolution_integral_030_050.pdf}
\sourcecaption{AMVF note Sec.\ 7.1, p.\ 14 shows the same kind of separation distribution; the fitted scale is our diagnostic. No truth-matching window enters this fit.}
\end{frame}

% Four independent examples of the agreed unweighted ntuple diagnostic.
% Every dot is one assigned track. The Split example includes both vertices;
% Fake is the unlabelled-track proxy, shown as a strict plurality below 50%.
\begin{frame}{Unweighted vertex categories: four examples}
\context{Each dot is one assigned track. Four independent examples of the convention used in our plots.}
\centering
\resizebox{.98\linewidth}{!}{%
\begin{tikzpicture}[x=1cm,y=1cm,font=\small]
  \colorlet{catA}{mydarkblue!85}
  \colorlet{catB}{myorange!85!black}
  \colorlet{catC}{teal!70!black}
  \colorlet{catF}{black!45}
  \foreach \x/\w in {1.6/3.0,5.0/3.0,9.35/5.0,13.7/3.0}{
    \node[pvcard,minimum width=\w cm,minimum height=3.65cm] at (\x,2.05) {};
  }
  % Clean: one truth interaction contributes 5/6 of the assigned tracks.
  \node[font=\small\bfseries,text=mygreen!65!black] at (1.6,3.5) {Clean / Matched};
  \foreach \p in {(1.0,2.75),(1.6,2.75),(2.2,2.75),(1.0,2.12),(1.6,2.12)}
    \fill[catA] \p circle (5pt);
  \fill[catB] (2.2,2.12) circle (5pt);
  \node[align=center,font=\footnotesize] at (1.6,1.1) {$5/6=83\%$ from one\\truth collision};
  % Merged: 60/40 avoids an ambiguous tie for the dominant truth origin.
  \node[font=\small\bfseries,text=myorange!85!black] at (5.0,3.5) {Merged};
  \foreach \p in {(4.4,2.75),(5.0,2.75),(5.6,2.75)}
    \fill[catA] \p circle (5pt);
  \foreach \p in {(4.7,2.12),(5.3,2.12)}
    \fill[catB] \p circle (5pt);
  \node[align=center,font=\footnotesize] at (5.0,1.1) {$60\%+40\%$\\Neither reaches $70\%$};
  % Split: both vertices are dominated by the same truth interaction.
  \node[font=\small\bfseries,text=myblue] at (9.35,3.5) {Split};
  \foreach \x/\name in {8.15/v_a,10.55/v_b}{
    \node[draw=black!25,rounded corners=3pt,minimum width=1.7cm,minimum height=1.1cm] at (\x,2.55) {};
    \node[font=\scriptsize,text=black!60] at (\x,3.08) {$\name$};
    \foreach \dx in {-.42,0,.42} \fill[catC] (\x+\dx,2.55) circle (5pt);
  }
  \node[font=\footnotesize] at (9.35,1.65) {$\sum p_T^2(v_a)>\sum p_T^2(v_b)$};
  \node[font=\footnotesize,text=mygreen!65!black] at (8.15,.98) {Keep Clean};
  \node[font=\footnotesize\bfseries,text=myblue] at (10.55,.98) {Relabel Split};
  \draw[-{Latex[length=2mm]},draw=myblue!65] (10.55,1.95)--(10.55,1.24);
  % Fake proxy: 3 unlabelled + 2 from each truth collision.
  \node[font=\small\bfseries,text=myred] at (13.7,3.5) {Fake (proxy)};
  \foreach \p in {(13.1,2.75),(13.7,2.75),(14.3,2.75)}
    \fill[catF] \p circle (5pt);
  \foreach \p in {(13.05,2.12),(13.48,2.12)} \fill[catA] \p circle (5pt);
  \foreach \p in {(13.92,2.12),(14.35,2.12)} \fill[catB] \p circle (5pt);
  \node[align=center,font=\footnotesize] at (13.7,1.1) {$3/7$ unlabelled $>2/7$\\from either collision};
  % Colours refer to truth origin within each independent example.
  \fill[catA] (1.0,-.45) circle (5pt);
  \node[anchor=west,font=\footnotesize] at (1.25,-.45) {truth collision A};
  \fill[catB] (4.7,-.45) circle (5pt);
  \node[anchor=west,font=\footnotesize] at (4.95,-.45) {truth collision B};
  \fill[catC] (8.4,-.45) circle (5pt);
  \node[anchor=west,font=\footnotesize] at (8.65,-.45) {same truth collision};
  \fill[catF] (12.65,-.45) circle (5pt);
  \node[anchor=west,font=\footnotesize] at (12.9,-.45) {unlabelled tracks};
\end{tikzpicture}}
\par\vspace{.65em}\small
Track counts set the category; $\sum p_T^2$ only resolves the Split label.
\sourcecaption{Unweighted approximation to \amvfnote. The paper uses fit weights and truth-classified fake tracks; neither is available here. Main results: top-1 GNN assignment, no score cut, reco $\ntrk\geq2$.}
\end{frame}


\begin{frame}{Vertex categories: our unweighted definitions}
\context{Every assigned track counts once. Truth labels identify its collision of origin.}
\small\renewcommand{\arraystretch}{1.5}
\begin{tabular}{p{3.2cm}p{10.3cm}}
\toprule
\textbf{Clean / Matched} & At least 70\% of the assigned tracks come from one true collision.\\
\textbf{Merged} & A true collision is the largest contribution, but none reaches 70\% of the assigned tracks.\\
\textbf{Split} & Several reco vertices have the same dominant true collision. Only the vertex with the largest $\sum p_T^2$ keeps Clean or Merged.\\
\textbf{Fake (proxy)} & The unlabelled-track group is the largest contribution.\\
\bottomrule
\end{tabular}
\par\vspace{.6em}\footnotesize
Split relabels duplicates after classifying their track content. Both methods retain reco vertices with at least two tracks. All category denominators include unlabelled tracks.
\sourcecaption{Adapted from \amvfnote; Track and Vertex Reconstruction with the ATLAS Inner Detector, Sec.\ 6.5, pp.\ 44--45 (local draft). Those papers use fit weights and truth-classified fake tracks. Our ntuples support this unweighted approximation.}
\end{frame}
